problem:  
Let $\Sigma_1, \Sigma_2 \subset S^3$ be two embedded, compact, nondegenerate minimal surfaces that intersect orthogonally along a closed geodesic $\gamma$. Assume that each $\Sigma_i$ is nondegenerate in the sense that its Jacobi operator has trivial kernel modulo the infinitesimal isometries of $S^3$.  

Consider a construction in which one replaces, along a small tubular neighborhood of $\gamma$, the local intersection region of $\Sigma_1 \cup \Sigma_2$ by a scaled Scherk-type minimal surface that desingularizes two orthogonally intersecting planes. For a collection of $k$ evenly spaced necks inserted along $\gamma$, denote by $\Sigma_{\varepsilon,k}$ the resulting surface depending on the neck size $\varepsilon>0$ and number $k\in\mathbb{N}$.  

**Question:**  
Under the above assumptions, does there exist a sequence $(\varepsilon_k,k)$ with $\varepsilon_k \to 0$ and $k \to \infty$ such that the corresponding embedded minimal surfaces $\Sigma_{\varepsilon_k,k}\subset S^3$ exist and:  
1. $\Sigma_{\varepsilon_k,k}$ are smooth, embedded, minimal surfaces of increasing genus, $\mathrm{genus}(\Sigma_{\varepsilon_k,k}) \to \infty$ as $k \to \infty$;  
2. $\Sigma_{\varepsilon_k,k}$ converge (as varifolds, or equivalently as stationary integral varifolds) to $\Sigma_1 \cup \Sigma_2$;  
3. the neck balancing conditions along $\gamma$ can be satisfied by suitable small perturbations preserving minimality and embeddedness?  

In other words, can every pair of orthogonally intersecting, nondegenerate compact minimal surfaces in $S^3$ be connected and desingularized along their common intersection, producing an infinite family of new, high-genus, embedded minimal surfaces in $S^3$ via Scherk-type gluing?  

Why is it a "good" problem:  
This problem sits precisely at the heart of modern geometric analysis. It asks whether one can universally carry out a stable desingularization of intersecting minimal surfaces in $S^3$ by Scherk-type gluing when both initial components are nondegenerate. The question is deceptively simple—can one smooth a clean geometric intersection—but its resolution would require a synthesis of deep analytical, geometric, and topological methods.  

1. **Profound Insight:**  
It explores whether the known gluing constructions (Kapouleas, Mazzeo–Pacard, Kapouleas–Yang) can be extended from symmetric examples to a broad, “universal” scheme in $S^3$. A positive result would suggest that minimal surfaces in the three-sphere form a connected moduli space dense in varifolds, revealing a structural completeness phenomenon.

2. **Bridge Across Fields:**  
The problem links nonlinear PDE theory (analysis of the Jacobi operator, perturbative existence), differential geometry (curvature and orthogonality in $S^3$), and geometric measure theory (varifold convergence and compactness). It bridges discrete geometric constructions with analytic deformation theory, connecting the geometry of individual examples with the topology of moduli spaces.

3. **Extreme Simplicity:**  
The statement is geometrically intuitive—whether two smooth minimal surfaces meeting orthogonally can be “glued smoothly"—but the techniques behind this are highly nontrivial. One must balance analytic constraints (kernel obstructions, neck-size matching) against global geometric topology (genus accumulation and embeddedness).

4. **Potential Impact:**  
A resolution would illuminate how complex the landscape of minimal surfaces in $S^3$ truly is, potentially opening paths toward a general theory classifying all embedded minimal surfaces by degenerations and desingularizations. It would also provide a new analytic framework for constructing families of minimal surfaces with controlled asymptotics and topology.

Thus, the problem represents a perfect blend of simplicity and depth, embodying the aesthetic and conceptual qualities of a "good" research problem.